% articulo-13-secciones.tex - Secciones y títulos.
% (needspace, titleformat).
%
% SECCIONES
%   \section siempre abre página nueva con drop folio (thispagestyle{plain}).
%   Variable YAML:
%     impresion: true  →  \section siempre en recto (impar), página en blanco
%                         intercalada sin encabezado ni folio si hace falta.
%     (sin declarar)   →  página nueva, recto o verso según caiga (por defecto).
%
%   \cleardoublepage: páginas en blanco intercaladas sin encabezado ni folio.
%   \cleardoubleoddpage: avanza a la siguiente página impar (recto). Requiere
%     que \c@page esté sincronizado con la paridad física.
%
% JERARQUÍA DE TÍTULOS
%   \section       16 pt, redonda
%   \subsection    12 pt, negrita
%   \subsubsection 12 pt, negrita cursiva
%   \paragraph     12 pt, cursiva         (en línea propia)
%   \subparagraph  12 pt, redonda         (en línea propia)
%
%   Principio tipográfico (Chicago): espacio generoso ANTES del título
%   (separa la sección anterior), mínimo DESPUÉS (el título y su texto
%   son una unidad visual).
%     \titlespacing*{cmd}{izq}{antes}{después}
%     El * suprime la sangría del párrafo inmediatamente posterior.
%
%   \titleformat*{} replica el estilo en la versión estrella (\section* etc.)
%   que pandoc emite cuando el título lleva {-} o {.unnumbered}.
%
%   Numeración: desactivada por defecto (numbersections: true para activar).
%   Para suprimir en un título concreto:
%     # Agradecimientos {-}
%     ## Subtítulo {.unnumbered}
%
% FAMILIA TIPOGRÁFICA DE LOS TÍTULOS
%   Variable YAML:
%     titulos-sans: true  →  todos los niveles de título en sans-serif (\sffamily).
%                            Convención habitual cuando el cuerpo va en serif y
%                            se busca contraste entre texto y aparato estructural.
%     (sin declarar)      →  títulos en serif, familia principal del documento
%                            (\rmfamily, por defecto).

% ── Secciones ─────────────────────────────────────────────────────────────────

\usepackage{needspace}

% Páginas en blanco intercaladas: sin encabezado ni folio.
\makeatletter
\renewcommand{\cleardoublepage}{%
  \clearpage
  \if@twoside
    \ifodd\c@page\else
      \hbox{}\thispagestyle{empty}\newpage
    \fi
  \fi
}
\makeatother

\let\oldsection\section

% \cleardoubleoddpage: avanza a la siguiente página impar (recto).
\makeatletter
\newcommand{\cleardoubleoddpage}{%
  \clearpage
  \ifodd\c@page\else
    \hbox{}\thispagestyle{empty}\clearpage
  \fi
}
\makeatother

% Redefinición de \section: drop folio en inicio de sección.
% En article, \section no llama a \thispagestyle por sí solo.
\renewcommand{\section}{%
  $if(impresion)$
  \cleardoubleoddpage
  $else$
  \clearpage
  $endif$
  \thispagestyle{plain}%
  \oldsection%
}

% ── Formato y espaciado de títulos ────────────────────────────────────────────

\titleformat{\section}
  {$if(titulos-sans)$\sffamily$else$\rmfamily$endif$\fontsize{16}{20}\selectfont}
  {\thesection}{1em}{}
\titleformat*{\section}
  {$if(titulos-sans)$\sffamily$else$\rmfamily$endif$\fontsize{16}{20}\selectfont}
\titlespacing*{\section}
  {0pt}{0pt}{2\baselineskip}

\titleformat{\subsection}
  {$if(titulos-sans)$\sffamily$else$\rmfamily$endif$\normalsize\bfseries}
  {\thesubsection}{1em}{}
\titleformat*{\subsection}
  {$if(titulos-sans)$\sffamily$else$\rmfamily$endif$\normalsize\bfseries}
\titlespacing*{\subsection}
  {0pt}{\baselineskip}{0pt}

\titleformat{\subsubsection}
  {$if(titulos-sans)$\sffamily$else$\rmfamily$endif$\normalsize\bfseries\itshape}
  {\thesubsubsection}{1em}{}
\titleformat*{\subsubsection}
  {$if(titulos-sans)$\sffamily$else$\rmfamily$endif$\normalsize\bfseries\itshape}
\titlespacing*{\subsubsection}
  {0pt}{\baselineskip}{0pt}

\titleformat{\paragraph}[hang]
  {$if(titulos-sans)$\sffamily$else$\rmfamily$endif$\normalsize\itshape}
  {\theparagraph}{1em}{}
\titleformat*{\paragraph}
  {$if(titulos-sans)$\sffamily$else$\rmfamily$endif$\normalsize\itshape}
\titlespacing*{\paragraph}
  {0pt}{\baselineskip}{0pt}

\titleformat{\subparagraph}[hang]
  {$if(titulos-sans)$\sffamily$else$\rmfamily$endif$\normalsize}
  {\thesubparagraph}{1em}{}
\titleformat*{\subparagraph}
  {$if(titulos-sans)$\sffamily$else$\rmfamily$endif$\normalsize}
\titlespacing*{\subparagraph}
  {0pt}{\baselineskip}{0pt}
